\section{Dokumentation des KI-Einsatzes}
\label{sec:ki-dokumentation}

Die folgende Tabelle dokumentiert den begrenzten unterstützenden Einsatz
generativer KI. Die fachlichen Inhalte und Konzeptionen der Arbeit wurden von
den Autoren selbst entwickelt.

\begin{center}
\footnotesize
\setlength{\tabcolsep}{3pt}
\renewcommand{\arraystretch}{1.05}
\begin{tabularx}{\textwidth}{|>{\raggedright\arraybackslash}p{3.1cm}|
  >{\raggedright\arraybackslash}X|
  >{\raggedright\arraybackslash}p{4.0cm}|
  >{\raggedright\arraybackslash}p{2.4cm}|}
  \hline
  \textbf{Anfrage bzw. Arbeitsschritt} &
  \textbf{Grund für die Verwendung der KI} &
  \textbf{Ergebnis und weitere Bearbeitung} &
  \textbf{Verwendete Software} \\
  \hline
  Zusätzliche Quellenrecherche &
  Unterstützung bei der Suche nach ergänzenden fachlich passenden
  Veröffentlichungen. &
  Vorschläge für zusätzliche Quellen; Auswahl, Prüfung und Zitation der
  tatsächlich verwendeten Quellen erfolgten durch die Autoren. &
  ChatGPT / Gemini \\
  \hline
  Recherche zum Unternehmens- und Anforderungskontext der Allianz &
  Suche nach ergänzenden Informationen zu den Geschäftsfeldern,
  Rahmenbedingungen, datenbezogenen Anforderungen sowie der Organisationsstruktur der Allianz. &
  Hinweise auf potenziell relevante Aspekte und Informationsquellen; die
  fachliche Einordnung und Übertragung auf das vorgegebene Szenario erfolgten
  durch die Autoren. &
  ChatGPT / Gemini \\
  \hline
  Benennung von Spalten im Datenschema &
  Unterstützung bei der Suche nach fachlich eindeutigen und einheitlichen
  Bezeichnungen für einzelne Attribute und Tabellenspalten. &
  Vorschläge für Spaltennamen. Aufbau, Inhalte, Tabellen, Beziehungen und die
  gesamte fachliche Logik des Datenschemas wurden von den Autoren selbst
  entwickelt. &
  ChatGPT / Gemini\\
  \hline
\end{tabularx}
\end{center}
